\documentclass[11pt,a4paper]{article}

\usepackage[a4paper,margin=25mm]{geometry}
\usepackage[T1]{fontenc}
\usepackage[utf8]{inputenc}
\usepackage{lmodern}
\usepackage{microtype}
\usepackage{hyperref}
\usepackage{longtable}
\usepackage{booktabs}
\usepackage{array}
\usepackage{listings}
\usepackage{xcolor}
\usepackage{enumitem}
\usepackage{tabularx}

\hypersetup{
  colorlinks=true,
  linkcolor=blue,
  citecolor=blue,
  urlcolor=blue
}

\lstdefinestyle{code}{
  basicstyle=\ttfamily\small,
  breaklines=true,
  columns=fullflexible,
  frame=single,
  backgroundcolor=\color{gray!5}
}

\lstset{style=code}

\title{Dedicated KVM Multi-Node Kubernetes Deployment of the Solana Kubernetes Observability Core}
\author{Oleksandr Khoshaba\\ORCID: 0000-0001-5375-6280}
\date{Version 0.4.0\\2026\\Technical report DOI: \href{https://doi.org/10.5281/zenodo.20561100}{10.5281/zenodo.20561100}}

\begin{document}

\maketitle

\begin{abstract}
This technical report documents the deployment and validation of the Solana Kubernetes Observability Core on a dedicated KVM-based multi-node Kubernetes cluster. The work extends the previous Minikube-based Kubernetes validation by deploying the same observability-oriented Solana workload on a kubeadm cluster composed of one control-plane virtual machine and two worker virtual machines. The validated deployment includes a Solana validator StatefulSet, validator ledger PersistentVolumeClaim, wallet-init Job, latency monitor Deployment, RPC Service, Yellowstone/Geyser gRPC Service, metrics Service, and wallet transfer validation Job. The stage remains focused on infrastructure and observability.
\end{abstract}

\tableofcontents

\section{Introduction}

The Solana Containerised Testbed is a reproducible local Solana environment for infrastructure, observability, dataset generation, and latency/performance research on legacy x86\_64 hardware. Earlier stages established the containerised no-AVX2 Solana local testbed~\cite{khoshaba_containerising_solana_v010}, added the Yellowstone/Geyser observability extension~\cite{khoshaba_yellowstone_geyser_v020}, and validated the Kubernetes Observability Core in a Minikube environment~\cite{khoshaba_solana_testbed_v030,khoshaba_kubernetes_observability_preprint_v030}.

This report documents the next stage: the deployment of the Solana Kubernetes Observability Core on a dedicated KVM-based multi-node Kubernetes cluster. The implementation and evidence for this stage are released as Solana Containerised Testbed v0.4.0~\cite{khoshaba_solana_testbed_v040}.

\section{Project Evolution}

The project evolved through several stages:

\begin{enumerate}[leftmargin=*]
  \item Compose-based local Solana testbed.
  \item No-AVX2 compatibility workflow.
  \item Yellowstone/Geyser observability extension.
  \item Minikube-based Kubernetes validation.
  \item Dedicated KVM multi-node Kubernetes deployment.
\end{enumerate}

\section{Host and Virtualisation Environment}

The deployment target is a CentOS Stream 9 host with KVM/libvirt support. VM disks are placed under the host home filesystem rather than the smaller root filesystem.

\section{Kubernetes Cluster Topology}

The validated cluster uses one control-plane VM and two worker VMs.

\begin{table}[htbp]
\centering
\small
\caption{Dedicated KVM Kubernetes cluster topology.}
\label{tab:kvm-cluster-topology}
\begin{tabularx}{\textwidth}{@{}p{0.22\textwidth}p{0.18\textwidth}p{0.27\textwidth}X@{}}
\toprule
Node & Role & Label & Main purpose \\
\midrule
\texttt{k8s-cp-01} &
control-plane &
\texttt{testbed-role= control-plane} &
Kubernetes API and control-plane services. \\

\texttt{k8s-worker-01} &
worker &
\texttt{testbed-role=validator} &
Solana validator workload and validator ledger PVC. \\

\texttt{k8s-worker-02} &
worker &
\texttt{testbed-role= observability} &
Monitor workload and auxiliary validation jobs. \\
\bottomrule
\end{tabularx}
\end{table}

\section{Kubernetes Bootstrap and Ansible Automation}

The infrastructure automation is stored under:

\begin{lstlisting}
infra/ansible/kvm-kubeadm/
\end{lstlisting}

The automation covers node preparation, containerd installation, Kubernetes tool installation, control-plane initialisation, worker join, CNI installation, node labelling, local-path storage installation, validation, and reset support.

\subsection{Main Ansible Entry Points}

The node preparation entry point combines the bootstrap, operating-system preparation, container runtime installation, Kubernetes tooling installation, and readiness validation steps.

\lstinputlisting{listings/site-prepare-nodes.yml}

The cluster creation entry point initialises the control plane, installs Calico, joins the workers, applies node labels, and installs the local-path provisioner.

\lstinputlisting{listings/site-create-cluster.yml}

The reset playbook documents the teardown path. This is important because a reproducible infrastructure baseline must support both deployment and reset.

\lstinputlisting{listings/reset-kubeadm-cluster.yml}

\section{Solana Kubernetes Observability Core}

The portable Kubernetes manifests are stored under:

\begin{lstlisting}
k8s/base/
\end{lstlisting}

The deployed core includes the validator StatefulSet, validator ledger PVC, wallet-init Job, monitor Deployment, RPC Service, Yellowstone/Geyser gRPC Service, and metrics Service.

\section{Kustomize Overlay for KVM Multi-Node Deployment}

The dedicated KVM multi-node overlay is stored under:

\begin{lstlisting}
k8s/overlays/kvm-multinode/
\end{lstlisting}

The overlay adds validator placement on \texttt{testbed-role=validator}, monitor placement on \texttt{testbed-role=observability}, and explicit \texttt{local-path} storage for the validator ledger PVC.

\section{Kustomize Overlay Listings}

The dedicated KVM multi-node deployment is represented by the \texttt{k8s/overlays/kvm-multinode} overlay. The overlay keeps the base manifests portable and adds only the environment-specific scheduling and storage rules.

\subsection{Overlay Kustomization}

\lstinputlisting{listings/kvm-multinode-kustomization.yaml}

\subsection{Validator Placement Patch}

\lstinputlisting{listings/patch-validator-placement.yaml}

\subsection{Monitor Placement Patch}

\lstinputlisting{listings/patch-monitor-placement.yaml}

\subsection{Validator Ledger PVC Patch}

\lstinputlisting{listings/patch-validator-ledger-pvc.yaml}

\section{Validation Evidence}

The validation evidence is stored in the repository under:

\begin{lstlisting}
results/kvm-kubeadm-cluster/
results/kvm-multinode-observability-core/
\end{lstlisting}

Selected evidence files were also copied into the report directory to make the technical report reproducible from the LaTeX source tree.

\subsection{Kubernetes Node Readiness and Placement Labels}

The cluster contains one control-plane node and two worker nodes. The worker nodes are labelled according to their intended testbed role.

\begin{lstlisting}
\lstinputlisting{evidence/nodes-labels.txt}
\end{lstlisting}

The validator workload is assigned to the node labelled \texttt{testbed-role=validator}. The observability workload is assigned to the node labelled \texttt{testbed-role=observability}. This explicit placement is the main difference from the earlier Minikube validation, where all workloads were scheduled onto a single-node cluster.

\subsection{StorageClass Validation}

The validator ledger PVC uses the \texttt{local-path} StorageClass. This is appropriate for the first dedicated KVM baseline because the validator ledger is intentionally node-local and bound to the validator worker.

\lstinputlisting{evidence/storageclass.txt}

\subsection{Workload Placement}

The deployed pods confirm the intended placement: the validator pod runs on \texttt{k8s-worker-01}, while the monitor runs on \texttt{k8s-worker-02}.

\lstinputlisting{evidence/pods-wide.txt}

\subsection{Services and EndpointSlices}

The deployment exposes three internal services: RPC on port \texttt{8899}, Yellowstone/Geyser gRPC on port \texttt{10000}, and metrics on port \texttt{9464}.

\lstinputlisting{evidence/services.txt}

EndpointSlice evidence confirms that the services resolve to the expected pod endpoints.

\lstinputlisting{evidence/endpointslices.txt}

\subsection{RPC Health}

The RPC health check returned \texttt{ok} through the Kubernetes Service endpoint.

\lstinputlisting{evidence/rpc-getHealth-clean.json}

\subsection{Wallet Transfer Validation}

The wallet transfer validation was executed as a Kubernetes Job. It created a temporary sender wallet, requested an airdrop, transferred \texttt{0.5 SOL} to the receiver, confirmed the transaction, and printed the final balances.

The validation Job name was recorded as follows:

\lstinputlisting{evidence/wallet-transfer-validation-job-name.txt}

The completed Kubernetes Job was recorded as follows:

\lstinputlisting{evidence/wallet-transfer-validation-job.txt}

The completed validation pod was recorded as follows:

\lstinputlisting{evidence/wallet-transfer-validation-pod.txt}

The validation confirmed that the sender balance changed from \texttt{2 SOL} to \texttt{1.499995 SOL}, while the receiver balance changed from \texttt{10000 SOL} to \texttt{10000.5 SOL}.

\subsection{Metrics After Wallet Transfer}

After the transfer validation, the metrics endpoint reported both slot interval observations and a transaction latency sample.

\lstinputlisting{evidence/metrics-after-wallet-transfer.txt}

The important line is:

\begin{lstlisting}
solana_transaction_latency_seconds_count 1
\end{lstlisting}

This confirms that the monitor observed at least one transaction latency sample in the dedicated KVM multi-node deployment.

\section{Scope and Limitations}

This report documents infrastructure and observability validation only. Load generation, dashboarding, Prometheus/Grafana integration, MPC, reinforcement learning, MARL, and Agave are out of scope.

\section{Conclusion}

The Solana Kubernetes Observability Core has been successfully deployed and validated on a dedicated KVM-based multi-node Kubernetes cluster.

\bibliographystyle{plain}
\bibliography{references}

\end{document}
